\documentclass[11pt]{article}
\usepackage[margin=1in]{geometry}
\usepackage{amsmath,amssymb,amsthm,booktabs,array}
\usepackage[T1]{fontenc}
\usepackage{lmodern}
\usepackage{microtype}
\usepackage[colorlinks=true,linkcolor=blue,citecolor=blue,urlcolor=blue]{hyperref}
\newtheorem{theorem}{Theorem}
\newtheorem{lemma}[theorem]{Lemma}
\newtheorem{proposition}[theorem]{Proposition}
\newtheorem{computation}[theorem]{Computation}
\newcommand{\SR}{\operatorname{SR}}
\newcommand{\wt}{\operatorname{wt}}
\newcommand{\PP}{P_{83}}
\title{An upper bound of 83 for the eighth Sidon--Ramsey number}
\author{Research draft for human review\\\small Automated research campaign, researcher 2}
\date{September 11, 2026}
\begin{document}
\maketitle
\begin{abstract}
A Sidon set of integers has distinct unordered pair sums, including sums
of an element with itself. The eighth Sidon--Ramsey number is the least
$n$ for which $\{1,\ldots,n\}$ cannot be partitioned into eight Sidon sets.
We give an exact computer-assisted proof that no such partition exists
for $n=83$, yielding $81\leq\SR(8)\leq83$. The proof covers every class-size
profile. An explicit integer weight vector reduces the problem to
weighted packings of eleven-element Sidon sets and complete searches in
their complements. The final profile has an exhaustive cover by
65,073,232 anchored triples. Two implementations agree on all case counts
and recorded nonempty completion queries. Public source, compact case
ledgers, and a complete replay command accompany the proof. The lower
bound uses a previously published partition of $\{1,\ldots,80\}$.
\end{abstract}

\section{Statement and context}
Write $[n]=\{1,\ldots,n\}$ and $P_n=\{0,\ldots,n-1\}$. A finite set
$A\subset\mathbb Z$ is \emph{Sidon} if
\[
 a+b=c+d,\quad a,b,c,d\in A
 \quad\Longrightarrow\quad \{a,b\}=\{c,d\}.
\]
Pairs here are unordered and repetitions are allowed. Translation and
reflection preserve this property. A Sidon $k$-partition is a partition
into at most $k$ Sidon sets, equivalently a coloring with $k$ colors whose
color classes are Sidon. We define $\SR(k)$ as the least $n$ with no such
partition of $[n]$, using the convention in \cite{EGMRS,EP}. For $n\geq k$,
requiring exactly $k$ nonempty classes is equivalent, since Sidon classes
can be split. The class-size bound below also rules out empty classes in
a hypothetical eight-coloring at $n=83$.

Espinosa-Garc\'ia and Pellicer established $\SR(8)\leq86$;
Theorem 4.1 of their open preprint gives the proof \cite{EP}. The
construction of Xiu, Fan and Liang \cite[Table I, $I=8$]{XFL} supplies
eight disjoint ten-element Sidon sets covering $[80]$. Our numerical
result is the following.

\begin{theorem}\label{thm:main}
There is no partition of $P_{83}$ into eight Sidon sets. Consequently,
\[
                         81\leq\SR(8)\leq83.
\]
\end{theorem}

The upper bound is established by the finite computations specified in
this paper and the accompanying source. We do not assume that a
hypothetical partition is balanced. In particular, eliminating only
profile $11^3 10^5$ would not suffice; the other six possible profiles
are handled first. Here $s^r$ means $r$ classes of size $s$.
The result does not determine whether $\SR(8)$ is 81, 82, or 83.

The proof has two computational stages. The first establishes the
complete profile landscape at $P_{83}$ and excludes every profile except
$11^3 10^5$. The second excludes that profile. Both are replayed from
fresh catalogs, and their interface is checked case by case.
Earlier exclusions at $P_{84}$ and $P_{85}$ are not mathematical premises.
The contribution is the finite numerical bound with its exhaustive
verification; no novelty is claimed for the elementary search reductions.

\section{Enumerating Sidon sets and certifying weights}\label{sec:weights}
For $a_1<\cdots<a_s$, the Sidon property is equivalent to distinctness of
all positive differences $a_j-a_i$, $i<j$. Indeed, equal positive
differences rearrange to equal unordered sums. Conversely, canceling a
common term in a nontrivial sum equality is impossible; arranging the
terms of such an equality gives two equal positive differences. This
also accounts for a collision involving a double, such as $a+c=2b$.

The two catalog methods use this equivalence. One extends increasing
prefixes while tracking used positive differences. The other fixes the
minimum and maximum elements and tracks unordered sums, including
doubles. Each tests every compatible next element, subject only to
necessary bounds: enough available points must remain; distinct positive
differences force diameter at least $\binom{s}{2}$; and $r$ future
consecutive gaps have sum at least $r(r+1)/2$. The difference method may
strengthen the last bound by using the $r$ smallest positive differences
not already used. Every completed set has a unique increasing order.
For normalized enumeration, all translates fitting the interval are
emitted. These observations justify completeness of the catalog
recursions; the implementation is part of the computational trust
boundary described in Section~\ref{sec:repro}.

For $D\subseteq\PP$, define $\wt(D)=\sum_{x\in D}u_x$, where the
nonnegative integer weights $u_x$ are given in
Table~\ref{tab:weights}. They satisfy $u_x=u_{82-x}$ and
\begin{equation}\label{eq:constants}
 W:=\wt(\PP)=31{,}134{,}774,\qquad
 M:=4{,}000{,}000,\qquad \max_x u_x=444{,}444.
\end{equation}
The vector is supplied explicitly. Its method of discovery is irrelevant
to validity: all needed inequalities are verified by integer enumeration.

\begin{table}[ht]
\centering
\begin{tabular}{rr@{\qquad}rr@{\qquad}rr}
\toprule
$x$&$u_x$&$x$&$u_x$&$x$&$u_x$\\
\midrule
0&243714&14&354517&28&428807\\
1&252661&15&361815&29&430468\\
2&254655&16&368123&30&434201\\
3&254974&17&375885&31&438601\\
4&260326&18&385230&32&440138\\
5&271212&19&394088&33&441273\\
6&285241&20&400351&34&443400\\
7&298729&21&405971&35&444362\\
8&307489&22&410588&36&443574\\
9&313975&23&412145&37&442251\\
10&321450&24&414224&38&442668\\
11&329981&25&417955&39&443991\\
12&337804&26&422030&40&444444\\
13&345652&27&426202&41&444444\\
\bottomrule
\end{tabular}

\caption{The 42 independent entries of the weight vector. Set
$u_{82-x}=u_x$ for the remaining entries.}\label{tab:weights}
\end{table}

\begin{computation}\label{comp:catalogs}
The complete catalogs in $\PP$ have the following sizes and weight maxima.
\[
\begin{array}{c|r|r}
 |A|&\text{number of Sidon sets}&\max\wt(A)\\\hline
 10&24{,}751{,}806&3{,}999{,}980\\
 11&15{,}958&3{,}999{,}979\\
 12&0&\text{not applicable}
\end{array}
\]
Both catalog methods give identical sets for sizes 10 and 11 and find
no size-12 set. The ten-element binary catalogs agree byte for byte;
the eleven-element catalogs agree as sets after sorting.
\end{computation}

A larger Sidon set would contain a size-12 subset, so every class has
size at most 11. Classes with at most nine points have weight at most
$9\cdot444{,}444=3{,}999{,}996<M$. Thus Computation~\ref{comp:catalogs}
proves the uniform bound
\begin{equation}\label{eq:cap}
                   \wt(A)\leq M\quad\text{for every Sidon }A\subseteq\PP.
\end{equation}
We deliberately retain this slightly loose common cap throughout the
proof, so all thresholds correspond exactly to the published case ledgers.

\begin{lemma}\label{lem:profiles}
Every Sidon eight-partition of $\PP$, if one exists, has one of these
seven profiles:
\begin{align*}
&11^7 6,\quad 11^6 10\,7,\quad 11^6 9\,8,\quad
 11^5 10^2 8,\quad 11^5 10\,9^2,\\
&11^4 10^3 9,\qquad 11^3 10^5.
\end{align*}
\end{lemma}
\begin{proof}
Writing each class size as $11-d_i$ gives nonnegative integers $d_i$
with $\sum_i d_i=88-83=5$. The seven integer partitions of 5 give the
listed profiles. This also shows that at least three classes have size 11.
\end{proof}

\section{Exhaustive packing and completion reductions}
We record the reductions that connect the finite computations to all
hypothetical partitions. For any choice of $r$ classes in an eight-class
partition, \eqref{eq:cap} implies
\begin{equation}\label{eq:union}
 \wt\left(\bigcup_{i=1}^{r}A_i\right)\geq W-(8-r)M.
\end{equation}
Disjointness makes the weight of the union additive.

\subsection{Reflection and canonical anchors}\label{sec:reflection}
Let $m(A)=\sum_{x\in A}2^x$ be the membership mask, and
$\rho(A)=\{82-x:x\in A\}$. No eleven-element Sidon set is fixed by
$\rho$: a fixed set of that size contains multiple distinct mirrored
pairs, each with sum 82. Partition the eleven-element catalog into
reflection pairs. In each pair choose the member with smaller mask as
representative. Sort the pairs by decreasing weight, breaking ties by
representative mask, and list each representative followed by its
reflection. Denote these rows by $A_0,\ldots,A_{15957}$.

\begin{lemma}\label{lem:anchor}
Every collection of $r\geq1$ disjoint eleven-element classes, up to
reflecting the entire collection, occurs among the increasing row
sequences
\[
        2q=i_1<i_2<\cdots<i_r,
\]
where $q$ is the earliest reflection pair represented in the collection.
If the collection can be part of an eight-class partition, then
\[
                    r\wt(A_{2q})\geq W-(8-r)M.
\]
\end{lemma}
\begin{proof}
If the earliest represented pair contributes its representative, use it
as the first row. Otherwise reflect the whole collection. Reflection
leaves the set of represented pair indices unchanged and interchanges
the two rows of the earliest pair. Every other selected row is later.
This remains valid if both members of that pair occur. Since row weights
are nonincreasing, the first weight bounds all $r$ selected weights.
Apply \eqref{eq:union}.
\end{proof}

This is a complete cover, not an assertion of one representative per
reflection orbit. Such multiplicities do not affect exclusion. After an
anchor is fixed, both packing traversals enumerate increasing disjoint
row sequences satisfying \eqref{eq:union}. One checks intersection with
the current union; the other filters compatible candidates recursively
(or intersects adjacency bitsets). Weight pruning uses upper bounds
from the heaviest still available rows. These bounds can discard only
sequences unable to reach \eqref{eq:union}.

\subsection{Heaviest-class recursion}\label{sec:recursion}
\begin{lemma}\label{lem:heaviest}
Suppose $D$ is to be partitioned into $k$ ten-element Sidon classes,
all of weight at most $U$. If $\wt(D)>kU$, no such partition exists.
Otherwise a heaviest class $T$ satisfies
\[
              \left\lceil\frac{\wt(D)}k\right\rceil
              \leq\wt(T)\leq U.
\]
Consequently, it suffices to enumerate every Sidon ten-set $T\subset D$
in this interval and recurse on $(D\setminus T,k-1,\wt(T))$.
For $k=1$, direct testing of $D$ decides the state.
\end{lemma}
\begin{proof}
The first statement follows by adding the $k$ class caps. The maximum
of $k$ weights is at least their average. Order the classes of any
partition by nonincreasing weight and apply this fact successively.
Every partition gives a branch of the recursion, including ties.
\end{proof}

For a domain to be partitioned into $t$ tens and one nine, cap the nine
by $M$. A heaviest ten then satisfies
\begin{equation}\label{eq:mixed}
 \left\lceil\frac{\wt(D)-M}{t}\right\rceil\leq\wt(T)\leq U,
 \qquad \wt(D)\leq tU+M.
\end{equation}
Recurse on the tens in nonincreasing order and test the final nine
directly. There is no requirement that the nine weigh less than the
last ten. This distinction is needed for coverage of the mixed profile.

The complete heavy-ten catalog consists of the tens satisfying
\begin{equation}\label{eq:cutoff}
       \wt(T)\geq C:=\left\lceil\frac{W-6M}{2}\right\rceil
                    =3{,}567{,}387.
\end{equation}
Filtering either complete ten catalog gives exactly 4,832,138 such sets.
In a balanced completion state with $k\geq2$, at most $8-k$ classes
have already been removed, so
\[
 \frac{\wt(D)}k\geq\frac{W-(8-k)M}{k}
             =M+\frac{W-8M}{k}\geq\frac{W-6M}{2}.
\]
The last inequality uses $W-8M=-865{,}226<0$ and $k\geq2$.
Thus the heavy catalog contains every ten needed by
Lemma~\ref{lem:heaviest}. The last ten is tested directly and need not
belong to the heavy catalog.

For the mixed profile, while $t\geq2$ tens remain together with the
nine, the same calculation gives
\[
 \frac{\wt(D)-M}{t}\geq\frac{W-(8-t)M}{t}
             \geq\frac{W-6M}{2}.
\]
When one ten and the nine remain, the domain has 19 points; every
Sidon ten in the interval \eqref{eq:mixed} is generated directly.
The heavy-catalog restriction is not imposed at that stage.

\section{Excluding the profiles with at least four elevens}
\subsection{Five or more elevens}
Any five eleven-element classes in a partition satisfy
\[
                       \wt(A_1\cup\cdots\cup A_5)\geq W-3M
                                                   =19{,}134{,}774.
\]
Their complement has 28 points. Lemma~\ref{lem:profiles} gives precisely
five possible remaining profiles:
\begin{equation}\label{eq:28profiles}
          (11,11,6),\ (11,10,7),\ (11,9,8),\ (10,10,8),\ (10,9,9).
\end{equation}
These include the possibility that the original partition has more than
five elevens.

\begin{computation}\label{comp:five}
There are exactly 2,142 unordered disjoint five-tuples satisfying the
weight threshold, with 2,142 distinct 28-point complements. None of these
complements admits any profile in \eqref{eq:28profiles}.
\end{computation}

The first packing method recursively filters compatible rows; the second
uses the full eleven-set disjointness graph. Both produce the same sorted
five-tuple list. Two complement routines generate Sidon subsets by
positive differences or unordered sums, respectively, and test all five
profiles. They find 66 ten-subset occurrences and no eleven-subset
occurrences across the domains. The recorded nine-subset count is 2:
this counts only the conditional nine queries actually required by the
search, not all nine-subsets of all 28-point domains.

For completeness, these query shortcuts follow from containment. Every
eleven-set contains eleven distinct ten-subsets, so a domain with fewer
than eleven tens cannot contain an eleven. Profiles with an eleven can
be decided by enumerating the relevant eleven/eleven, eleven/ten, or
eleven/nine disjoint choices and testing the last class. With no eleven,
only $(10,10,8)$ and $(10,9,9)$ remain. Enumerate their ten choices and,
for the latter, the required nines in the residual 18-point domain;
test the complement directly. Positive controls contain an explicit
partition of every profile in \eqref{eq:28profiles}.

\subsection{Exactly four elevens}
The only profile is $11^4 10^3 9$. The union of its four elevens has
weight at least $W-4M=15,134,774$. Lemma~\ref{lem:anchor} gives 2,701
eligible anchors; all increasing disjoint anchored quadruples above
this threshold are traversed.

\begin{computation}\label{comp:four}
The cover contains 7,805,097 anchored quadruples. None of their 39-point
complements partitions into three Sidon tens and one Sidon nine.
\end{computation}

The completion procedure is \eqref{eq:mixed}. Two subset-query engines
supply every heavy ten contained in the current domain: one uses a
radix tree with common-point masks and maximum weights, the other incidence
bitsets over a catalog ordered by weight. The former discards a subtree
only when a common point is forbidden or its maximum weight is below
the query minimum; leaves also enforce the query maximum, while the latter intersects the weight range with the complements
of the incidence sets for forbidden points. Thus both implement the
same exact subset query. Their complete recorded nonempty query streams
agree byte for byte. The recursion totals are:
\begin{center}
\begin{tabular}{lrr}
\toprule
Remaining profile&States tested&Candidate tens returned\\
\midrule
$10^3 9$&7,805,097&7,813,293\\
$10^2 9$&7,813,293&1,077,835\\
$10\,9$&1,077,835&380\\
$9$&380&---\\
\bottomrule
\end{tabular}
\end{center}
Each of the 380 terminal nine-sets has an explicit collision of unordered
sums, independently checked in Python. The per-anchor ledger records
all 2,701 anchors, including those without a qualifying quadruple.

\section{Excluding the final profile}\label{sec:balanced}
Only $11^3 10^5$ remains. The three elevens must have total weight at
least
\begin{equation}\label{eq:triple}
                        W-5M=11{,}134{,}774.
\end{equation}
The canonical anchor condition is $3\wt(A_{2q})\geq11,134,774$.
The ordered catalog gives exactly 5,364 eligible anchors. Two preliminary
triple enumerations, one using direct intersections and the other bitset
adjacency, establish the following cover.

\begin{computation}\label{comp:triples}
The complete anchored cover of triples satisfying \eqref{eq:triple}
contains 65,073,232 triples. Of the 5,364 anchors, 207 have no such triple
and 5,157 have at least one. The two per-anchor ledgers agree exactly.
\end{computation}

For each triple, apply Lemma~\ref{lem:heaviest} to its 50-point complement,
starting with $k=5$ and $U=M$. The heavy catalog is sufficient by
\eqref{eq:cutoff}. Every branch ends either at a necessary infeasibility
bound, an empty subset query, or a directly tested residual ten.

\begin{computation}\label{comp:balanced}
Neither of the two complete packing-and-completion traversals finds a
partition for any of the 65,073,232 anchored triples. Their per-anchor
counts, recorded nonempty query streams, and terminal records agree.
\end{computation}

The resulting search counts are shown in Table~\ref{tab:balanced}. There
are 1,771 terminal occurrences, each a distinct failing residual ten,
associated with 1,644 distinct anchored triples. Each terminal record
lists the three eleven-row indices, four chosen tens, the residual ten,
and an explicit pair-sum collision. Python checks every listed class,
disjointness, coverage of $P_{83}$, weight ordering, and the collision.

\begin{table}[ht]
\centering
\begin{tabular}{rrr}
\toprule
Tens remaining&States tested&Candidate tens returned\\
\midrule
5&65,073,232&185,585,643\\
4&185,585,643&66,521,099\\
3&66,521,099&2,261,478\\
2&2,261,478&1,771\\
1&1,771&0\\
\bottomrule
\end{tabular}
\caption{Complete recursion for profile $11^3 10^5$. Counts are traversal
occurrences, not necessarily distinct residual domains.}\label{tab:balanced}
\end{table}

An additional check removes the weight-order restriction from the final
two classes. For every terminal record, take the union of its last chosen
ten with the failing residual ten. There are 1,754 distinct 20-point
domains. Two direct enumerators find 4,040 ten-subset occurrences across
these domains and no partition of any domain into two Sidon tens.
This checks the terminal decisions independently of the heavy-catalog
query. It does not replace the global traversal that reaches them.

\begin{proof}[Proof of Theorem~\ref{thm:main}]
Computation~\ref{comp:catalogs} and Lemma~\ref{lem:profiles} cover all
possible profiles. Computation~\ref{comp:five} excludes the first five;
Computation~\ref{comp:four} excludes the sixth. Lemmas~\ref{lem:anchor}
and~\ref{lem:heaviest} ensure that the remaining profile is covered by
Computations~\ref{comp:triples} and~\ref{comp:balanced}, which exclude it.
Translation gives the same exclusion on $[83]$. The partition of $[80]$
in Appendix~\ref{app:80} gives the lower bound.
\end{proof}

\section{Reproduction and scope of verification}\label{sec:repro}
The source is available in the \texttt{sidon\_ramsey\_8} directory of
\cite{source}. The consolidated \texttt{paper} directory contains this
manuscript, a source/input hash manifest, and \texttt{verify.py}.
From the repository root, the complete command is
\begin{verbatim}
python3 sidon_ramsey_8/paper/verify.py --work /tmp/sr8-full-proof
\end{verbatim}
The work directory must be new and outside the repository. The command
runs \path{p83_profiles/reproduce.py} and
\path{p83_exclusion/reproduce.py}, including sanitizer controls, then
checks the interface. It regenerates the full ten- and eleven-catalogs
in both stages; compares them byte for byte between stages; and matches
all 5,364 triple-cover rows to the completion ledger. The published
expected counts and hashes are consistency checks, not substitutes for
complete enumeration. A separate aggregation mode checks existing
completed runs and labels that narrower action explicitly.

The two stages have no imported P84 exclusion. The profile driver reads
an earlier 84-entry weight vector only to check the provenance identity
$u_i=v_i+v_{i+1}$; it verifies all P83 weight bounds afresh. A general
catalog program located under \texttt{p84\_weight\_certificates} is run
on the explicitly supplied domain $P_{83}$. The manifest records both
of these source dependencies, avoiding dependence on an undocumented
working copy or a prior search checkpoint.

The profile stage includes positive examples of all five 28-point
profiles and mixed ten/nine completions. The balanced stage compares
4,712 small exact-cover decisions with a Python definition-level
solver, checks positive 50-point and 20-point domains from the known
partition, and rejects a malformed duplicated catalog. AddressSanitizer
and UndefinedBehaviorSanitizer run on the small and positive controls.
The full production traversals use optimized builds; sanitizer runs
cover the controls, not the entire production search.

For the balanced stage the two nonempty query streams have
8,044,688,696 bytes each; for the four-eleven stage they have 299,368,488
bytes each. Each stream is split into shards, compared with its counterpart,
and checked against case counts and completion markers. Empty queries
are counted by the state ledgers but are not individually recorded in
these binary streams. Compact ledgers and representative terminal
records are committed; bulky catalogs and traces are regenerated in the
external work directory. About 25 GB of free disk space is a practical
allowance for the complete default replay.

All mathematical quantities in the proof computations are integers.
Point and difference masks fit in 128-bit unsigned words; pair-sum masks
use two words when needed. The supported parameter ranges keep every
shift below its word width, and counts use 64-bit integers. Floating-point
arithmetic is used for elapsed-time reporting, not exclusion decisions.
Python optimization must be disabled because the stage drivers use
assertions as proof checks. The trust boundary includes the coverage
arguments, source implementations, compiler, runtime, and hardware.
There is no proof-assistant formalization or portable SAT-style
unsatisfiability certificate.

A separate automated research identity replayed both stages from fresh
inputs and issued an acceptance with limitations \cite{review}. Its
additional unordered-sum checker enumerated all nine-subsets of the
2,142 full 28-point domains (318,212 occurrences), rather than only the
conditional nine queries used by the producer, and separately checked
the terminal 20-point domains. It did not supply a third independently
written global traversal. The two global implementations here were
developed within the same research campaign, and agreement between them
is a computational cross-check rather than independent authorship.
The review is not journal peer review. This manuscript remains a draft
for human mathematical review and authorship assignment.

\appendix
\section{The known lower-bound partition}\label{app:80}
The following rows are the $I=8$ construction in
\cite[Table I]{XFL}, reproduced as finite mathematical data. Their union
is $[80]$, they are disjoint, and each row has 55 distinct unordered
pair sums including its ten doubles. These properties are checked
directly by the replay entry point.
\begin{center}\small
\begin{tabular}{rrrrrrrrrr}
18&20&32&39&47&50&63&67&72&73\\
11&12&17&21&34&37&45&52&64&66\\
8&10&19&25&26&46&49&54&68&80\\
6&9&16&27&35&40&60&62&76&77\\
5&7&15&22&28&31&56&70&74&75\\
4&13&36&38&42&43&55&58&69&79\\
3&14&23&29&30&48&51&53&61&65\\
1&2&24&33&41&44&57&59&71&78
\end{tabular}
\end{center}

\section{Compact artifact identities}
The profile stage's three-eleven cover ledger has SHA-256
\begin{center}\small\texttt{d9858a448b47d98b3a3879e7a0f14e0f773d}\\
\texttt{c50519fb25c8af635830569af236}\end{center}
and the completion stage's ledger has SHA-256
\begin{center}\small\texttt{b32d9a00e42e5349dbf6677988fb7c9da107}\\
\texttt{fd325b5507cd308392006a57b250}.\end{center}
Each displayed pair of lines is one concatenated digest. The first
ledger records each anchor weight and triple count; the second records
the same triple count and completion-call totals. These differing file
formats explain why their digests differ despite exact agreement at
the interface. Full digests for source files, weights, catalogs, and
terminal data appear in the machine-readable manifests and replay
reports. A digest identifies an artifact; it does not itself prove
completeness.

\begingroup
\raggedright
\bibliographystyle{plain}
\bibliography{references}
\endgroup
\end{document}
